\renewcommand{\arraystretch}{1.3}
\captionsetup[table]{justification=raggedright,singlelinecheck=false}
\begin{table}[H]
\caption{Temuan wawancara praktisi sebagai bukti empiris permasalahan}
\label{tbl:bukti_wawancara}
\centering
\begin{tabular}{| p{0.18\textwidth} | p{0.42\textwidth} | p{0.25\textwidth} |}
\hline
\textbf{Tema Keluhan} & \textbf{Temuan Konkret dari Praktisi} & \textbf{Frekuensi/Intensitas} \\
\hline
Inefisiensi penelusuran (B01) &
Kesulitan menelusuri relasi antarobjek dan mencari konteks pada berkas konfigurasi panjang; penjelasan kerap tidak tuntas. &
Penggunaan harian; sebagian mencapai 15--20 kueri per hari \\
\hline
Halusinasi dan ketidakandalan (B02) &
Jawaban LLM gagal dijalankan, berbeda antarmodel, serta membuat asumsi yang tidak sesuai konteks pengguna. &
Dilaporkan seluruh narasumber \\
\hline
Risiko konfigurasi (B03) &
Keluaran konfigurasi YAML tidak konsisten antarcara bertanya; konfigurasi tidak memperbarui variabel terbaru. &
Berulang pada skenario generatif \\
\hline
\end{tabular}
\end{table}
